\begin{abstract}
Recent advances in Large Language Models (LLMs) offer unique opportunities for scientific tasks, yet their ability to reason over complex numerical data remains largely unexplored. While several works investigate the ability of LLMs to create textual descriptions of scientific objects (e.g., proteins), we ask whether they can understand those objects and answer comparative questions about them. This kind of reasoning beyond description is crucial for the development of useful co-scientists. We test two different astrophysical datasets, involving both stars and galaxies, with pre-trained multimodal features that serve as the LLM data tokens. Adding those tokens as additional information, we fine-tune an 8B Llama model on various scientific reasoning tasks. We find that while the model can describe the stellar objects reasonably well, they fail when asked about populations, outlier detections, and cross-modal correlations. We further show that by training a small non-language head that would highlight the scientific information (e.g., a flow matching model for population questions), and adding its latent features as an extra token, the performance of the LLM on the relevant tasks improves dramatically. The results are consistent across two independent datasets and highlight the lack of statistical reasoning abilities in current language-based models and the need for dedicated models to transform the data in a way that helps LLMs understand it. 
\end{abstract}